\subsection{Magnus 代数}

设 \(\hat{\mathbf{A}}(\mathbf{X})\) 是积模 \(\prod_{n\geqslant 0}\mathbf{A}^n (\mathbf{X})\) 。我们在 \(\hat{\mathbf{A}}(\mathbf{X})\) 上按规则

\[
(a. b) _ {n} = \sum_ {i = 0} ^ {n} a _ {i}. b _ {n - i}\tag{1}
\]

其中 \(a = (a_{n})\) 与 \(b = (b_{n})\) 都属于 \(\hat{\mathbf{A}}(\mathbf{X})\) 。我们知道 (《交换代数》， 第 III 章， § 2， no. 12， 例 1) \(\hat{\mathbf{A}}(\mathbf{X})\) 是一个结合代数，并且 \(\mathbf{A}(\mathbf{X})\) 与 \(\hat{\mathbf{A}}(\mathbf{X})\) 的由除有限多项外全为零的序列组成的子代数等同。

给 \(\hat{A}(X)\) 赋予诸因子上离散拓扑的乘积拓扑

\(A^{n}(X)\) ；当环 K 取离散拓扑时，这个拓扑使 \(\hat{A}(X)\) 成为一个完备 Hausdorff 拓扑代数，并且 \(A(X)\) 在 \(\hat{A}(X)\) 中稠密。设 \(a = (a_{n}) \in \hat{A}(X)\) ；族 \((a_{n})_{n \geqslant 0}\) 是可和的且 \(a = \sum_{n \geqslant 0} a_{n}\) 。

对每个整数 \(m \geqslant 0\) ，用 \(\hat{\mathbf{A}}_m(\mathbf{X})\) 表示由级数 \(a = \sum_{n \geqslant m} a_n\) 组成的理想，这些级数满足 \(a_n \in \mathrm{A}^n(\mathbf{X})\) （对一切 \(n \geqslant m\) ）。这个理想序列是 \(\hat{\mathbf{A}}(\mathbf{X})\) 中 0 的一个邻域基本系，并且是 \(\hat{\mathbf{A}}(\mathbf{X})\) 上的一个整滤过。与上述滤过相伴的阶函数记作 \(\omega\) ；那么 \(\omega(0) = +\infty\) ，并且 \(\omega(a) = m\) （若 \(a = \sum_{n \geqslant m} a_n\) ，其中 \(a_n \in \mathrm{A}^n(\mathbf{X})\) （对一切 \(n \geqslant m\) ）且 \(a_m \neq 0\) ）(§ 4， no. 1 与 no. 2)。

\(\hat{\mathbf{A}}(\mathbf{X})\) 称为集合 \(\mathbf{X}\) 的 Magnus 代数（系数在 \(\mathbf{K}\) 中）。若 \(\mathbf{K}\) 有歧义，我们记作 \(\hat{\mathbf{A}}_{\mathbb{R}}(\mathbf{X})\) 。

命题 1. 设 B 是一个带实滤过 \((\mathrm{B}_{\alpha})_{\alpha \in \mathbb{R}}\) 的幺结合代数，使得 B 是 Hausdorff 的且完备的 (§ 4， no. 1 与 no. 2)。设 \(f\) 是 \(\mathbf{X}\) 到 B 中的一个映射，使得存在 \(\lambda > 0\) 满足 \(f(\mathbf{X}) \subset \mathbf{B}_{\lambda}\) 。那么 \(f\) 可以以唯一的方式扩张为一个连续幺同态 \(f\) （从 \(\hat{\mathbf{A}}(\mathbf{X})\) 到 B 中）。

设 \(f'\) 是 A(X) 到 B 中的扩张 \(f\) 的唯一的幺代数同态 (《代数》， 第 III 章， \S 2， no. 7， 命题 7)。我们证明 \(f'\) 是连续的： \(f'(A^n(X)) \subset B_{n\lambda}\) ，从而 \(f'(\hat{A}_n(X) \cap A(X)) \subset B_{n\lambda}\) 。因此 \(f'\) 可以以唯一的方式由连续性扩张为一个同态 \(\hat{A}(\mathbf{X}) \to \mathbf{B}\) 。

我们保留命题 1 的假设与记号，并设 \(u \in \hat{\mathbf{A}}(\mathbf{X})\) 。元素 \(f(u)\) 记作 \(u((f(x))_{x \in \mathbf{X}})\) ，称为把 \(f(x)\) 代换 \(x\) （在 \(u\) 中）所得的结果。特别地， \(u((x)_{x \in \mathbf{X}}) = u\) 。现设 \(\boldsymbol{u} = (u_y)_{y \in \mathbf{Y}}\) 是 \(\hat{\mathbf{A}}(\mathbf{X})\) 的元素的一个族，并设 \(v \in \hat{\mathbf{A}}(\mathbf{Y})\) 。由上所述我们可以定义元素 \(v((u_y)_{y \in \mathbf{Y}}) \in \hat{\mathbf{A}}(\mathbf{X})\) 。它记作 \(v \circ \boldsymbol{u}\) 。由于 \(u_y((f(x))) \in \mathbf{B}_{\lambda}\) ，可以把元素 \(u_y((f(x)))\) 代换 \(y\) （在 \(v\) 中）。于是映射 \(v \mapsto (v \circ \boldsymbol{u})((f(x))_{x \in \mathbf{X}})\) 与 \(v \mapsto v((u_y((f(x)))_{x \in \mathbf{X}})_{y \in \mathbf{Y}})\) 是 \(\hat{\mathbf{A}}(\mathbf{X})\) 到 \(\mathbf{B}\) 中的两个连续幺代数同态，它们取相同的值 \(u_y((f(x)))\) （在元素 \(y \in \mathbf{Y}\) 处）。因此 (命题 1)

\[
(v \circ \boldsymbol {u}) ((f (x))_{x \in \mathbf{X}}) = v ((u _ {y} ((f (x)))_{x \in \mathbf{X}})_{y \in \mathbf{Y}})\tag{2}
\]

对一切 \(v \in \hat{\mathrm{A}}(\mathrm{Y})\) 成立。

